\begin{afterword}
\summaryhead

The post introduces Gillian Russell's essay ``Epistemic Viciousness in the Martial Arts''. Russell tells of believing, at eleven, a teacher's promise that three years of punching the air would make it possible to kill a bull with one blow. The author summarizes the essay's account of how such beliefs arise: the art and the teacher are treated as sacred; practitioners have invested years in their techniques; beginners cannot learn from a book and must trust a teacher; historical masters are deferred to; and there are too few data to test techniques, or the methods of training them. Worst, martial artists live in this poverty of data but act as if they could afford to believe what they are told.

The author adds a memory that turned out not to be in the essay: martial arts declined once real fights stopped showing which schools were better, and videos now show high-ranked black belts, and a master who believed in ki, beaten by real fighters. Every one of these risk factors, the post says, carries over to any attempt at a ``rationality dojo''. It asks what can be done.

\responsehead

Most of the post is a summary of Russell's essay, and an accurate one. It follows the essay's sections in order, and its quotations match the draft on Russell's website closely. The author also corrects, openly, a memory of the essay that turned out to be wrong.

The post's own additions rest on memory. The history of decline after real fights ended has no source, and the videos offered for it show black belts losing today, not that things were once better. The same history supports ``A Sense That More Is Possible'', posted the same day. The ki master, recalled without a link, is probably Yanagi Ryuken, whose filmed defeat a commenter linked two days later. The recollection itself is borne out. And Russell's essay makes a point close to the one the author remembered: judo, where a throw either works or does not, leaves ``relatively little room for pretence'', while odd beliefs do better in arts with less competitive sparring.

The post's own claim is its last paragraph: ``Every single one of these risk factors transfers straight over to any attempt to start a `rationality dojo'.'' The post asserts this without going through the factors. For data poverty the claim is one the author argues elsewhere at length. For the need to trust a teacher, Russell's reason is that a martial art ``can't be learned from a book'', and the post does not say whether that holds for an art the author has so far taught in writing.

The post ends by asking, ``What can be done about it?'' For a reader who wants a start, the essay it summarizes has partial answers that the summary leaves out: test parts of a technique when the whole cannot be tested, and keep justified deference to a teacher apart from unjustified deference.

In short: a faithful summary of Russell's essay, with a history of decline added from memory and no source, and a claim that every risk factor carries over to rationality training that the post asserts without examining.
\end{afterword}
